\subsection{Model selection}
Hyperparameter screening used a deterministic 10\% sample from 2017--2020
and validation on 2021. Candidate configurations were then refitted on the
2017--2021 training pool and compared on a deterministic 10\% sample of 2022.
The grid covered LR $\texttt{regParam}=\{0.001,0.01,0.1,1\}$, RF
$(\texttt{numTrees},\texttt{maxDepth})=\{(50,6),(100,6),(50,8)\}$ and GBT
$(\texttt{maxDepth},\texttt{maxIter})=\{3,4\}\times\{20,40\}$; 2021 mean
six-threshold PR-AUC selected one setting per family, with mean Brier as the
tie-breaker.
The model-selection training sample contained 2,742,964 base windows and the
2022 evaluation sample contained 711,112. The five candidates were
climatology, LR-A, RF-B, GBT-A and GBT-B. Mean PR-AUC across the six primary
thresholds ranked candidates, with mean Brier score as the tie-breaker
(Table~\ref{tab:selection}).

The climatology baseline estimates training-data exceedance rates by station,
month, hour and threshold, shrinking each level toward the broader rate with
$\lambda=100$:
\[
  \hat p_{0,\tau}=\frac{X_{0,\tau}}{N_0},\qquad
  \hat p_{m,h,\tau}=\frac{X_{m,h,\tau}+100\hat p_{0,\tau}}{N_{m,h}+100},\qquad
  \hat p_{s,m,h,\tau}=\frac{X_{s,m,h,\tau}+100\hat p_{m,h,\tau}}{N_{s,m,h}+100}.
\]
Here $X$ and $N$ are positive and total training counts; an empty group takes
its parent rate, so unseen station-month-hour groups fall back to month-hour
and then overall rates.

\begin{table}[htbp]
\centering\small
\caption{Mean metrics across the six primary thresholds on the 2022 model-
selection sample. Each candidate used the same 2,742,964 sampled training
windows and 711,112 evaluation windows. Fit time is model-fitting time.}
\label{tab:selection}
\begin{tabular}{l r r r}
\toprule
Candidate & Mean PR-AUC & Mean Brier & Fit time (s) \\
\midrule
Climatology & 0.03847 & 0.02261 & --- \\
LR-A & 0.03253 & 0.02265 & 46.6 \\
RF-B & 0.50399 & 0.01466 & 255.9 \\
GBT-A & 0.03553 & 0.02286 & 319.7 \\
GBT-B & \textbf{0.52207} & \textbf{0.01402} & 383.5 \\
\bottomrule
\end{tabular}
\end{table}

Feature set A uses station, cyclic time and threshold. Feature set B adds
current and historical same-station rainfall, missingness and window-support
summaries. GBT-B used depth 4 and 40 iterations. The 2022 sample was used for
model selection; final generalisation performance is reported on the 2024
holdout.

\begin{table}[htbp]
\centering\small
\caption{GBT feature-set ablation on the deterministic 2022 model-selection
sample. Both variants used 2,742,964 training and 711,112 evaluation base
windows.}
\label{tab:feature-ablation}
\begin{tabular}{l l r r r}
\toprule
Variant & Added inputs over Set A & Mean PR-AUC & Mean Brier & Fit time (s) \\
\midrule
GBT-A & None & 0.03553 & 0.02286 & 319.7 \\
GBT-B & Same-station rainfall history & \textbf{0.52207} &
\textbf{0.01402} & 383.5 \\
\bottomrule
\end{tabular}
\end{table}

On the 2022 selection sample, GBT-B outperformed GBT-A, so Set B was retained
for the final model.

